% sec6.tex -- writer W4. Straight runs and design T.
\section{Straight runs and design T}\label{sec:designT}

Design T consists of two straight runs. On a straight run the true candidates obey two exact integrals, and so do the first and
second derivatives of the front in the label; this section collects these exact facts (valid for every $v\in(0,1)$) and defines
design T.

\emph{Heading frame.} While $P$ moves with constant unit velocity (its \emph{heading}), we use coordinates in which $P$ is at the
origin and the heading is $1\in\C$; a candidate is at $r=\rho e^{i\beta}=z+iy$, and for $0<\beta<\pi$ we put $\tau:=\tan(\beta/2)$.

\begin{proposition}\label{prop:runs}
On a straight run, a true candidate with $\rho>0$ satisfies
\[
\rho'=v-\cos\beta,\qquad\rho\,\beta'=\sin\beta;\qquad h:=\rho(\cos\beta+v)\ \text{ satisfies }\ h'=-(1-v^2).
\]
If $0<\beta<\pi$, then $\beta$ stays in $(0,\pi)$ and $C:=\rho\,\sg(\beta)$ is constant, where
\[
\sg(\beta):=\frac{\sin\beta}{\tan^v(\beta/2)}=\frac{2\tau^{1-v}}{1+\tau^2},\qquad\frac{\sg'}{\sg}=\frac{\cos\beta-v}{\sin\beta}.
\]
The function $\sk:=(\cos\beta+v)/\sg$ is strictly decreasing on $(0,\pi)$, from $+\infty$ to $-\infty$; $\sg$ is increasing on
$(0,\arccos v]$ and decreasing on $[\arccos v,\pi)$, and
\[
1=\sg(\pi/2)\le\gmax:=\sg(\arccos v)=(1+v)^{(1+v)/2}(1-v)^{(1-v)/2}\le e^{v^2}.
\]
\end{proposition}
\begin{proof}
$r'=v\rhat-1$; the components along $\rhat$ and $\rhat^\perp$ give the first two equations. Then
$h'=(\cos\beta+v)(v-\cos\beta)-\rho\sin\beta\,\beta'=v^2-\cos^2\beta-\sin^2\beta$. Since $\beta'=0$ at $\beta\in\{0,\pi\}$, $\beta$
cannot leave $(0,\pi)$. On $(0,\pi)$, $\log\sg=\log\sin\beta-v\log\tan(\beta/2)$ gives $\sg'/\sg=\cot\beta-v/\sin\beta$, and, as $\beta'/\sin\beta=1/\rho$,
$\rho'/\rho=(v-\cos\beta)\beta'/\sin\beta=-(\sg'/\sg)(\beta)\,\beta'$; so $C$ is constant. With $\sin\beta=2\tau/(1+\tau^2)$ we get the $\tau$-form of $\sg$, so $\sg\to0$ at
both ends of $(0,\pi)$. Next,
$\sk'=[-\sin\beta-(\cos\beta+v)(\cos\beta-v)/\sin\beta]/\sg=-(1-v^2)/(\sg\sin\beta)<0$, and $\cos\beta+v\to1+v$, resp.\ $-(1-v)$,
at $0$, resp.\ $\pi$, which gives the limits of $\sk$. The monotonicity of $\sg$ follows from the sign of $\cos\beta-v$. At
$\cos\beta=v$: $\sin\beta=(1-v^2)^{1/2}$ and $\tan(\beta/2)=((1-v)/(1+v))^{1/2}$, which gives the closed form. Finally
$\log\gmax=\frac12[(1+v)\log(1+v)+(1-v)\log(1-v)]=\sum_{n\ge1}\frac{v^{2n}}{2n(2n-1)}\le v^2\sum_{n\ge1}\frac1{2n(2n-1)}=v^2\log2\le v^2$.
\end{proof}

For a $1$-Lipschitz path $P$ and a time $t$, the \emph{true front} is the curve $\theta\mapsto S_\theta(t)$ on the labels whose true
candidates exist on $[0,t]$. The field $x\mapsto(x-P(t))/|x-P(t)|$ is $C^\infty$ in $x\ne P(t)$ and continuous in $t$; hence near
every label whose candidate has had $\rho>0$ on $[0,t]$ the front is $C^\infty$ in $\theta$, and $q:=\partial_\theta S_\theta$
satisfies $\partial_tq=v(\Pi/\rho)q$, $\Pi=I-\rhat\rhat^{\mathsf T}$. We write $w=|q|$, $T=q/w=e^{i\psi}$, $\phi=\arg r-\psi$ and
$k=\partial_s\psi=\partial_\theta\psi/w$ (so $k=1$ for the unit circle), as for virtual fronts (Definition~\ref{def:flow}).

\begin{lemma}\label{lem:truefront}
At every label and time with $\rho>0$ on $[0,t]$, the true front satisfies
\[
\frac{w_t}w=\frac{v\sin^2\!\phi}\rho,\qquad\psi_t=-\frac v\rho\sin\phi\cos\phi,\qquad
k_t=\frac v\rho(1-3\sin^2\!\phi)\,k+\frac v{\rho^2}\sin\phi\,(3\cos^2\!\phi-1),
\]
and $w\ge1$. If the true or the virtual candidate $\theta$ has $\rho\ge\eps$ on $[0,t]$, then $x_\theta=S_\theta$ on $[0,t]$; if this
holds for all labels of an open set, the true and the virtual front coincide there with all their label derivatives.
\end{lemma}
\begin{proof}
Let $n:=\rhat^\perp$; then $T=\cos\phi\,\rhat-\sin\phi\,n$, $\Pi T=-\sin\phi\,n$, $T^\perp=\cos\phi\,n+\sin\phi\,\rhat$. From
$\partial_t(wT)=v\Pi(wT)/\rho$: $w_t=vw\,T\cdot\Pi T/\rho=vw\sin^2\!\phi/\rho$, and
\[
T_t=\frac v\rho(\Pi T-(T\cdot\Pi T)T)=-\frac v\rho[\sin\phi\cos^2\!\phi\,n+\sin^2\!\phi\cos\phi\,\rhat]=-\frac v\rho\sin\phi\cos\phi\,T^\perp .
\]
At $t=0$, $w=|\partial_\theta e^{i\theta}|=1$, and $w$ is non-decreasing. For the curvature: $\partial_s\rho=\cos\phi$,
$\partial_s\rhat=\Pi T/\rho=-(\sin\phi/\rho)n$, so $\partial_s\phi=-\sin\phi/\rho-k$. Since $k=\partial_\theta\psi/w$ and the front is
smooth, $k_t=\partial_s(\psi_t)-k\,w_t/w$, and
\[
\partial_s\Bigl(-\frac v{2\rho}\sin2\phi\Bigr)=\frac v{\rho^2}\sin\phi\cos^2\!\phi+\frac v\rho\cos2\phi\Bigl(k+\frac{\sin\phi}\rho\Bigr);
\]
so $k_t=\frac v\rho k(\cos2\phi-\sin^2\!\phi)+\frac v{\rho^2}\sin\phi(\cos^2\!\phi+\cos2\phi)$, which is the claim. Finally, $F(x)=x/|x|$
for $|x|\ge\eps$ (Lemma~\ref{lem:field}(i)): a true candidate with $\rho\ge\eps$ on $[0,t]$ solves the virtual equation, whose
solutions are unique (Lemma~\ref{lem:flow}(a)); a virtual candidate with $\rho\ge\eps$ solves the true equation, whose solutions are
unique while $\rho>0$.
\end{proof}

From now on in this section, $P$ moves with constant unit velocity on a time interval, coordinates are those of the heading frame,
and candidates have $0<\beta<\pi$.

\begin{lemma}\label{lem:coord}
\begin{enumerate}
\item[(a)] $\rho=\frac C2(\tau^{v-1}+\tau^{v+1})$, $y=\rho\sin\beta=C\tau^v$, $z=\rho\cos\beta=\frac C2(\tau^{v-1}-\tau^{v+1})$, and
$h=\frac C2[(1+v)\tau^{v-1}-(1-v)\tau^{1+v}]$.
\item[(b)] $\tau\mapsto(1+v)\tau^{v-1}-(1-v)\tau^{1+v}$ is a strictly decreasing bijection $(0,\infty)\to\R$. Hence $\tau$ is strictly
increasing in time, and the state at any time of the run is determined by $(C,h)$.
\item[(c)] $\min_\tau\rho=C/\gmax$, attained at $\tau^2=(1-v)/(1+v)$, i.e.\ $\cos\beta=v$; $\rho$ is decreasing in time while
$\cos\beta>v$.
\item[(d)] If $h>0$: $\tau\le((1+v)C/(2h))^{1/(1-v)}$. If $h<0$: $\tau\ge(2|h|/((1-v)C))^{1/(1+v)}$, and
\[
\tau^{1+v}=\frac{2|h|}{(1-v)C}\Bigl[1+\frac{(1+v)C\tau^{v-1}}{2|h|}\Bigr].
\]
\end{enumerate}
\end{lemma}
\begin{proof}
(a) $\sin\beta=2\tau/(1+\tau^2)$, $\cos\beta=(1-\tau^2)/(1+\tau^2)$, $\rho=C/\sg(\beta)$ and $h=z+v\rho$. (b) The derivative is
$-(1-v^2)(\tau^{v-2}+\tau^v)<0$, the limits are $\pm\infty$, and $h'<0$. (c) $\rho'=v-\cos\beta$ (Proposition~\ref{prop:runs}), and
$\frac d{d\tau}(\tau^{v-1}+\tau^{v+1})=\tau^{v-2}((1+v)\tau^2-(1-v))$. (d) By (a), $(1+v)\tau^{v-1}=2h/C+(1-v)\tau^{1+v}\ge2h/C$, resp.\
$(1-v)\tau^{1+v}=2|h|/C+(1+v)\tau^{v-1}\ge2|h|/C$.
\end{proof}

\begin{lemma}[transport]\label{lem:transport}
Let $\theta\mapsto S_\theta$ be a $C^2$ family of true candidates with $0<\beta<\pi$ on the run, $q:=\partial_\theta S_\theta$
(heading frame), $w=|q|$, $q=we^{i\psi}$, $\phi=\beta-\psi$, and
\[
\Xi:=q+v\,e^{i\beta}\,\bar q .
\]
\begin{enumerate}
\item[(a)] $\operatorname{Re}\Xi=\partial_\theta h=:\sa$ and $\operatorname{Im}\Xi=\tau^v\partial_\theta C=:\tau^v\sfc$, where
$\sa(\theta)$, $\sfc(\theta)$ do not depend on time; $q=(\Xi-ve^{i\beta}\bar\Xi)/(1-v^2)$ and $(1-v)w\le|\Xi|\le(1+v)w$.
\item[(b)] Let $Y_r:=\sfc\tau^v/\sa$ (where $\sa\ne0$), $N_\psi:=(1+v\cos\beta)Y_r-v\sin\beta$, $D_\psi:=1-v\cos\beta-v\sin\beta\,Y_r$.
Then $(1-v^2)q=\sa(D_\psi+iN_\psi)$. If $\sa>0$ and $D_\psi>0$: $|\psi|<\pi/2$, $\tan\psi=N_\psi/D_\psi$, $|\psi|\le|\tan\psi|$ and
$w=\sa|D_\psi+iN_\psi|/(1-v^2)$.
\item[(c)] $J:=\sa\,\partial_\theta\sfc-\sfc\,\partial_\theta\sa$ does not depend on time, and
\[
k=\frac{\tau^vJ}{(1-v^2)\,w^3}-\TS,\qquad
\TS:=\frac{v\sin\phi}{(1-v^2)\,\rho}\Bigl[\frac{\rho\,\sa\,\partial_\theta\log C}{w^2}-v-\cos(\phi-\psi)\Bigr].
\]
\item[(d)] For two times $t'<t''$ of the run (data $'$, $''$): $\tau''\ge\tau'$,
\[
k''=\Bigl(\frac{\tau''}{\tau'}\Bigr)^v\Bigl(\frac{w'}{w''}\Bigr)^3(k'+\TS')-\TS'',\qquad\text{so}\qquad
|k''|\le\Bigl(\frac{\tau''}{\tau'}\Bigr)^v(|k'|+|\TS'|)+|\TS''|,
\]
and $w''\le(\tau''/\tau')^v|\Xi'|/(1-v)\le\frac{1+v}{1-v}(\tau''/\tau')^vw'$.
\item[(e)] At every time, $|\TS|\le\dfrac{v|\sin\phi|}{(1-v^2)\rho}\Bigl[\dfrac{\rho|\sa||\partial_\theta\log C|}{w^2}+1+v\Bigr]$
and $|\TS|\le\dfrac{v(3+v)}{2(1-v)y}$.
\end{enumerate}
\end{lemma}
\begin{proof}
(a) Write $q=X+iY$. As $P$ does not depend on $\theta$: $\partial_\theta z=X$, $\partial_\theta y=Y$,
$\partial_\theta\rho=X\cos\beta+Y\sin\beta$, $\rho\,\partial_\theta\beta=-X\sin\beta+Y\cos\beta=-w\sin\phi$. Since
$ve^{i\beta}\bar q=v(X\cos\beta+Y\sin\beta)+iv(X\sin\beta-Y\cos\beta)$,
\[
\partial_\theta h=X+v\,\partial_\theta\rho=\operatorname{Re}\Xi .
\]
By Lemma~\ref{lem:coord}(a), $C=y\tau^{-v}$, and $\partial_\theta\log\tau=\partial_\theta\beta/\sin\beta$; so
$\tau^v\partial_\theta C=Y-v\rho\,\partial_\theta\beta=Y(1-v\cos\beta)+vX\sin\beta=\operatorname{Im}\Xi$. Since $C$ is constant and
$h+(1-v^2)t$ is constant (Proposition~\ref{prop:runs}), $\sa$ and $\sfc$ do not depend on time. The map
$L:x\mapsto x+vR_\beta x$, $R_\beta x:=e^{i\beta}\bar x$, is $I+vR_\beta$ with $R_\beta$ a reflection; so $L^{-1}=(I-vR_\beta)/(1-v^2)$
and the singular values of $L$ are $1\pm v$.

(b) With $\Xi=\sa(1+iY_r)$: $(1-v^2)q=\Xi-ve^{i\beta}\bar\Xi$ and
$e^{i\beta}(1-iY_r)=\cos\beta+Y_r\sin\beta+i(\sin\beta-Y_r\cos\beta)$ give $(1-v^2)q=\sa(D_\psi+iN_\psi)$. If $\sa,D_\psi>0$, then
$\operatorname{Re}q>0$, so $|\psi|<\pi/2$ and $|\psi|\le|\tan\psi|$.

(c) $J$ is built from the time-independent $\sa,\sfc$. Put $V:=\Xi=\sa+i\sfc\tau^v=Lq$. Then
$\partial_\theta V=\partial_\theta\sa+i(\tau^v\partial_\theta\sfc+v\sfc\tau^v\partial_\theta\beta/\sin\beta)$,
$\partial_\beta(L^{-1})=-vR'_\beta/(1-v^2)$ with $R'_\beta x:=ie^{i\beta}\bar x$, and
$\partial_\theta q=L^{-1}\partial_\theta V+\partial_\theta\beta\,\partial_\beta(L^{-1})V$. For a real $2\times2$ matrix $M$,
$(Ma)\times(Mb)=\det M\,(a\times b)$, and $\det L^{-1}=1/(1-v^2)$; hence
\[
kw^3=q\times\partial_\theta q=\frac{V\times\partial_\theta V}{1-v^2}-\frac{v\,\partial_\theta\beta}{1-v^2}\,q\times R'_\beta(Lq).
\]
Here $V\times\partial_\theta V=\tau^vJ+v\sa\sfc\tau^v\partial_\theta\beta/\sin\beta$ and $\sfc\tau^v/\sin\beta=\rho\,\partial_\theta\log C$
(Lemma~\ref{lem:coord}(a): $\tau^v/\sin\beta=\tau^{v-1}(1+\tau^2)/2=\rho/C$). Further $Lq=w(e^{i\psi}+ve^{i\phi})$,
$R'_\beta(Lq)=iw(e^{i\phi}+ve^{i\psi})$, $q\times R'_\beta(Lq)=w^2(\cos(\phi-\psi)+v)$, and $\partial_\theta\beta=-w\sin\phi/\rho$.
Inserting gives
$kw^3=\frac{\tau^vJ}{1-v^2}-\frac{vw\sin\phi}{(1-v^2)\rho}[\rho\sa\,\partial_\theta\log C-w^2(v+\cos(\phi-\psi))]$.

(d) Eliminate $J$ between the formulas of (c) at $t'$ and $t''$. $\tau$ increases in time (Lemma~\ref{lem:coord}(b)) and $w$ is
non-decreasing (Lemma~\ref{lem:truefront}), so $w'/w''\le1$. $\Xi''=\sa+i\tau''^v\sfc$ gives $|\Xi''|\le(\tau''/\tau')^v|\Xi'|$; then use
(a).

(e) The first bound is (c) with $|v+\cos(\phi-\psi)|\le1+v$. For the second: $\partial_\theta\log C=\sfc/C=\operatorname{Im}\Xi/y$ and
$\sa=\operatorname{Re}\Xi$, so $\rho|\sa\,\partial_\theta\log C|/w^2\le\rho|\Xi|^2/(2yw^2)\le(1+v)^2\rho/(2y)$; with $|\sin\phi|\le1$ and
$\rho\ge y$: $|\TS|\le\frac v{1-v}[\frac{1+v}{2y}+\frac1\rho]\le\frac{v(3+v)}{2(1-v)y}$.
\end{proof}

\begin{definition}[design T]\label{def:T}
Let $t_1:=1/(1-v)+1$ $(=2+\mu)$.
\begin{itemize}
\item \emph{Run 1.} $P(t):=t$ for $0\le t\le t_1$.
\item \emph{Run 2.} Let $\alpha_0:=\arccos(v+1/t_1)$ and write $S_{\alpha_0}(t_1)-P(t_1)=\rho_0e^{i\beta^*}$ ($\rho_0>0$), $\nhat:=e^{i\beta^*}$.
Then $P(t):=P(t_1)+(t-t_1)\nhat$ for $t_1\le t\le\ttd$, where
\[
\ttd:=t_1+\frac{\rho_0-\eps}{1-v}+\eps\int_0^1\frac{d\sigma}{1-v\sigma(2-\sigma)} .
\]
\end{itemize}
Further $\ta:=t_1+(\rho_0-\eps)/(1-v)$, $\tc:=t_1+\rho_0/(1-v)$, $X_0:=P(\ttd)$ and $\omega(t):=\tc-t$.
\end{definition}

Here $v+1/t_1=v+(1-v)/(2-v)\in(\frac12,1)$, and $\rho_0>0$ because in run 1 the candidate $\alpha_0\in(0,\pi/2)$ has
$C=\sg(\alpha_0)>0$, so $\rho\ge C/\gmax$ (Lemma~\ref{lem:coord}(c)); $\rho_0>\eps$ by Lemma~\ref{lem:tang}(b), so $t_1<\ta<\ttd<\tc$.
$P$ is $1$-Lipschitz, $P(0)=0$, and $P$ depends only on $(v,\eps)$: $\nhat$ is computed from the state of the true candidate $\alpha_0$
at $t_1$, which is determined by run 1.

\emph{Run-1 data.} For $\alpha\in(0,\pi)$ let $(\rho_1,\beta_1)=(\rho_1(\alpha),\beta_1(\alpha))$ be the polar coordinates of
$S_\alpha(t_1)-P(t_1)$ in the heading frame of run 1, $\tau_1:=\tan(\beta_1/2)$, $\tau_\alpha:=\tan(\alpha/2)$ and
$y(\alpha,t_1):=\operatorname{Im}S_\alpha(t_1)$. In run 1 the lower half is the mirror image of the upper one,
$S_{2\pi-\alpha}=\overline{S_\alpha}$. At $t=0$, $\rho=1$, $\beta=\alpha$; hence, in run 1,
\begin{gather*}
C=\sg(\alpha),\qquad h(t)=\cos\alpha+v-(1-v^2)t,\qquad h_1:=h(t_1)=-(1-v^2)-(1-\cos\alpha)<0,\\
\sa=-\sin\alpha,\qquad\sfc=(\cos\alpha-v)\tau_\alpha^{-v},
\end{gather*}
using $(1-v^2)t_1=(1+v)(2-v)$ and $\sg'=\sg(\cos\alpha-v)/\sin\alpha$.

\emph{Run-2 data.} For $\theta\in(0,2\pi)\setminus\{\alpha_0\}$ let $(\sr(\theta),\sbb(\theta))$ be the polar coordinates of
$S_\theta(t_1)-P(t_1)$ about $\nhat$ (the heading frame of run 2), $\tau_{\sbb}:=\tan(\sbb/2)$, and, for run 2,
\[
C:=\sr\,\sg(\sbb),\qquad\sa:=\partial_\theta[\sr(\cos\sbb+v)],\qquad\sfc:=\partial_\theta C,\qquad\sA:=\partial_\theta\sr .
\]
From now on $C$, $h$, $\sa$, $\sfc$ refer to run 2 unless said otherwise. The \emph{window} is $I_\eta:=\{\theta:|\theta-\alpha_0|\le\eta\}$,
$\eta=0.11639$ (Section~\ref{sec:touch}), and $\xi:=|\theta-\alpha_0|$. In the frame of run 2 at time $t$ we write $S_\theta(t)-P(t)=(z+iy)\nhat=\rho e^{i\beta}\nhat$
and $e^{i\psi}\nhat$ for the unit tangent. Since $(1-v^2)(\tc-t_1)=(1+v)\rho_0$, every label whose true candidate has $\rho>0$ on run 2 (including $\alpha_0$, for $t<\tc$) has in run 2
\begin{equation}\label{eq:h2}
h(t)=\sr(\cos\sbb+v)-(1-v^2)(t-t_1)=(1+v)(\sr-\rho_0)-\sr(1-\cos\sbb)+(1-v^2)\,\omega(t).
\end{equation}
Whenever $h(t)>0$, Lemma~\ref{lem:coord}(d) gives $\tau\le((1+v)C/(2h))^{1+\mu}$, hence $y=C\tau^v\le C^{1+\mu}((1+v)/(2h))^\mu$, and
$C\le2\sr\tau_{\sbb}^{1-v}$ gives $C^{1+\mu}\le(2\sr)^{1+\mu}\tau_{\sbb}$; therefore
\begin{equation}\label{eq:tau2}
\tau\le\tau_{\sbb}\Bigl(\frac{(1+v)\sr}h\Bigr)^{1+\mu},\qquad y\le2\sr\,\tau_{\sbb}\Bigl(\frac{(1+v)\sr}h\Bigr)^{\mu}\qquad(h>0).
\end{equation}
